% ============================================================
% Chapter: Price recursion as the rational theory of reward
% Compiled from:
%   marketgrad-theory/Price recursion is the rational theory of reward.md   (backbone)
%   marketgrad-theory/Bregman divergences, mirror ascent and reward as
%       hyperstitious information.md                                        (Section \ref{agency:sec:hyperstition})
%   marketgrad/README.md, marketgrad/SPEC.md                                (Section \ref{agency:sec:library})
% ============================================================

\chapter{Price recursion as the rational theory of reward}\label{ch:agency}

\chaptersource{This chapter is based on the informal articles
``Reward
is hyperstitional information'' \citep{pallavisudhirRewardHyperstitional2026} and ``Price recursion is the rational theory of reward'' \citep{pallavisudhirPriceRecursion2026}
(originally titled ``Markets are equivalent to stuff''), written as part of the author's work
at MATS~10.0 in Richard Ngo's stream, together with the accompanying \texttt{marketgrad}
library (Section~\ref{agency:sec:library}) developed there.}

\chapterabstract{\draft{Rational behaviour for a purely predictive agent is completely
specified by Bayes's theorem; but for any agent that acts in the world there is an arbitrary
thing called the ``utility function'', shaped by an equally arbitrary reward. This chapter
develops an economic answer to where reward comes from. First, entropic mirror ascent ---
the natural optimization dynamic on probability space --- is shown to be continuous-time
Bayesian inference, so that reward is (hyperstitional) information. Then, in a general model
of a production economy, shadow prices are shown to obey a recursion that is simultaneously
backpropagation and Bellman recursion: chain economies are exactly multi-layer perceptrons,
markets are exactly MDPs with matrix-valued discount factors, and entropic mirror ascent on
routing shares is exactly replicator dynamics/continuous Bayesian inference, with revenue as
the evidence rate. The economy is thus the ``super-agent'' that generates reward for its
sub-agents, and price recursion is the rational theory of reward. The accompanying
\texttt{marketgrad} library implements this theory and numerically verifies its
identities.}}


Several
papers discuss RL algorithms based on sequentially auctioning a unit ``right to act'' among
agents---most recently Caspar Oesterheld's BRIA \citep{oesterheldTheoryBoundedInductive2023}
but also Baum's Hayek
Machine \citep{baumModelIntelligenceEconomy1999}, its predecessors the classifier systems
\citep{hollandPropertiesBucketBrigade1985a, schmidhuberEvolutionaryPrinciplesSelfreferential1987, schmidhuberLocalLearningAlgorithm1989} and
its successors Hayek-for-POMDPs \citep{kweeMarketBasedReinforcementLearning2001} and
Hayek with Vickrey auctions \citep{changDecentralizedReinforcementLearning2020}. In
``Market-based architectures in RL and beyond'' \citep{pallavisudhirMarketbasedArchitectures2025}
I suggested that a better market model for RL might come out of factoring the ``right to
act'' into ``goods'' axes, just like in real markets.

In this chapter we will formalize this intuition, and explain that in fact, markets are
equivalent to a lot of things: backprop, RL, and continuous Bayesian inference (aka
``replicators'').

\begin{draftblock}
The chapter is organized as follows. Section~\ref{agency:sec:hyperstition} reproduces the
precedent article ``Reward is hyperstitional information'', which develops the machinery of
Bregman divergences and entropic mirror ascent and argues that reward, under this dynamics,
is exactly (hyperstitional) Bayesian evidence; the market dynamics of
Section~\ref{agency:sec:cbi} instantiate this. Sections~\ref{agency:sec:setup}
to~\ref{agency:sec:rationalize} are the main article: the economic model and price
recursion, the market--MDP and market--MLP equivalences, continuous Bayesian inference as
market dynamics, and the resulting rationalization of reward.
Section~\ref{agency:sec:library} briefly describes the \texttt{marketgrad} library, and
Appendix~\ref{agency:sec:caveats} collects the main article's caveats and open questions.
\end{draftblock}

% ============================================================
\section{Bregman divergences, mirror ascent and reward as hyperstitious information}
\label{agency:sec:hyperstition}
% Source: marketgrad-theory/Bregman divergences, mirror ascent and reward as
% hyperstitious information.md (the precedent article, lifted in full).

The logarithmic scoring rule $b\log q(x)$ exactly
puts a price on information by the usual measure \citep{pallavisudhirLMSRSubsidy2024}: the expected score an agent with some
belief distribution $q$ gets is exactly $-bH(p,q)$ where $p$ is the true distribution (or,
more accurately: the score we expect that agent to get is $-bH(p,q)$, from our
perspective/belief $p$).

Consequently, the value of some information $Y$ is precisely its mutual information
$bI(X;Y)$ with the event of interest, etc.

Here's a question: \textbf{does every proper scoring rule correspond to an information
measure?} For instance, does the quadratic scoring rule incentivize obtaining information in
the sense of ``quadratic negentropy'' much like logarithmic scoring incentivizes obtaining
information in the sense of ``logarithmic negentropy''?

\subsection{Generalized entropies and Bregman divergence}\label{agency:sec:bregman}

A scoring rule is some function $S(x, q)$ that tells you how much an agent gets for belief
$q$ when an outcome $x$ is realized---equivalently that is how someone with belief
$\delta_x$ would score belief $q$. This also tells you how someone with belief $p$ would
score belief $q$:
$$G(p,q)=\sum_x p(x)S(x,q)$$
Properness means that $q=p$ maximizes it, i.e.
$$G(p):=G(p,p)=\max_q G(p,q)$$
$G(p)$ is the score we give to our own beliefs, thus there is no bias penalty, only an
uncertainty penalty: it measures our ``certainty'' about our own beliefs.

(So $G(p)$ and $G(p,q)$ are the generalized ``negentropy'' and ``cross-negentropy''
respectively, and their negatives are the generalized entropy and cross-entropy.)

Now observe: fix $q$ and look at the maps $p\mapsto G(p)$ and $p\mapsto G(p,q)$. The latter
is linear in $p$, and they touch at the point $p=q$ (everywhere else $G(p,q)$ lies below
$G(p)$).

Thus for every $q$, $G(p,q)$ is the tangent plane under the graph of $G(p)$ at the point
$p=q$. In other words, \textbf{any (neg-)entropy function $G(p)$ naturally induces a cross
(neg-)entropy function $G(p,q)$} by its tangent planes:
$$G(p,q)=G(q)+\langle\nabla G(q), p-q\rangle$$

When $G(p)$ is convex, i.e.\ it sits above all its tangent planes (in fact it is exactly the
max of all its tangent planes), the resulting scoring rule is proper i.e.\
$G(p)=\max_q G(p,q)$. \textbf{Thus every convex function generates an information theory.}

So the Generalized Scoring $S(x,q)=G(\delta_x, q)$, the Generalized Cross-Entropy $-G(p,q)$
and the Generalized Entropy $-G(p)$ are all equally fundamental; you can derive the rest
from any one of them.

We also define the Generalized KL-divergence (``Bregman Divergence'') as the vertical gap
between $G(p)$ and the tangent plane $G(p,q)$, i.e.\ ``how much does $p$ believe itself to
be better than $q$'':
$$D(p\parallel q)=G(p)-G(p,q)$$

For log score $S(x, q)=\log q(x)$, we have usual information theory. For the quadratic score
$S(x,q)=-\|q-\delta_x\|^2$ we have cross-negentropy $G(p,q)=\langle p, q\rangle-\|q\|^2-1$,
negentropy $G(p)=\|p\|^2-1$ (the negative of which, the quadratic entropy, is called the
``Gini impurity'') and Bregman divergence $D(p\parallel q)=\|p-q\|^2$.

An entropy function inducing a cross-entropy function may remind you of how norms naturally
induce inner products (polarization identity). So is cross-entropy an inner product? Not
really. If you take a quadratic form $G(p)=\frac12 \langle p, A p\rangle$ (differing from
the above quadratic entropy by constant term and factor) then $\nabla H(p)=Ap$, and
$$H(p,q)= \frac12 \langle q, A q\rangle+\langle Aq, p-q\rangle=\langle p,Aq\rangle-\frac12 \langle q,Aq\rangle$$
Which is not the inner product, but it does suggest a generalized ``inner product'' of
distributions: $\langle p, q\rangle =G(p,q)+G(q)=G(p)+G(q)-D(p\parallel q)$. But it is not
in general symmetric, etc.

\subsection{Mirror ascent}\label{agency:sec:mirror}

What exactly is gradient ascent?
\begin{quote}
``Climbing in the direction of steepest ascent''
\end{quote}
ok, what is ``the direction of steepest ascent?''
\begin{quote}
Among all $\varepsilon$-steps in each direction, we choose the $\varepsilon$-step that gives
the highest rise in the objective (as measured by a first-order Taylor approximation of the
objective, i.e.\ the gradient)
\end{quote}

I.e.\ given a current value $x_t$, we want to find the $x$ that maximizes $f(x)-f(x_t)$. The
first-order Taylor approximation for this is $\langle\nabla f(x_t), x-x_t\rangle$---but we
cannot trust the first-order Taylor approximation if we move too far, so we restrict the
search to an $\varepsilon$-neighbourhood around $x_t$. Thus:
$$x_{t+1}=\argmax_x \langle\nabla f(x_t), x-x_t\rangle, \quad D(x, x_t)\le\varepsilon$$

This is a constrained optimization problem, so define a Lagrangian:
$$\mathcal{L}(x,\lambda)= \langle\nabla f(x_t), x-x_t\rangle - \lambda (D(x, x_t)-\varepsilon)$$

Differentiating we get:
\begin{equation}
\nabla D(x_{t+1}, x_t) = \frac1\lambda \nabla f(x_t) \tag{$*$}
\end{equation}
$$D(x_{t+1}, x_t)=\varepsilon$$

For the usual Euclidean metric, $\nabla D(x_{t+1}, x_t) = 2(x_{t+1}-x_t)$ so we can solve
and get usual gradient ascent:
$$x_{t+1} = x_t + \eta \nabla f(x_t)$$

Instead suppose $D(x_{t+1}, x_t)$ is an arbitrary Bregman divergence, i.e.\
$D(x, y) = G(x) - G(y) - \langle \nabla G(y), x-y\rangle$ for convex $G$. The gradient of
this with respect to its first parameter is simply
$$\nabla_x D(x, y) = \nabla G(x) - \nabla G(y)$$

So plugging this back into $(*)$ we get the update rule:
\begin{equation}
\nabla G(x_{t+1}) = \nabla G(x_t) + \eta\nabla f(x_t)\tag{$**$}
\end{equation}

For convex $G$, there is a one-to-one map between the values of $x$ and the values of
$\nabla G(x)$; i.e.\ the gradients live in a ``$G$-dual space'' or \textbf{mirror space} and
this amounts exactly to doing gradient ascent in the mirror space, i.e.\ ``mirror ascent''.

For $G(x)= \|x\|^2$ (the usual quadratic metric), the map is the identity (scaled by $2$),
and---

\subsection{Reward is hyperstitional information}\label{agency:sec:hyperstitional}

for $G(p)=p\cdot\log(p)$ we have $\nabla G(p) = 1+\log(p)$ thus the update rule is
multiplicative (\textbf{entropic mirror ascent}):
$$p_{t+1} \propto p_t\exp\left(\eta\nabla f(p_t)\right)$$

(we are allowed to normalize because for probabilities the original problem is itself a
constrained optimization problem with a simplex constraint)

In other words: \textbf{optimizing some objective $f(p)$ over probability space with
entropic mirror ascent, is continuous-time Bayesian inference with likelihood function
$\exp\left(\eta\nabla f(p)\right)$}.

Unlike usual Bayesian inference:
\begin{enumerate}
  \item this likelihood function depends on the current $p$
  \item this likelihood function can be \ldots\ anything.
\end{enumerate}

From a predictive processing/active inference perspective, all rational behaviour must be
Bayesian inference, and utility functions are fake. What you see as ``utility
functions''/goals are actually the agent's prior: \textbf{and inference arbitrages between
the agent's beliefs and the environment} (when the environment's information updates the
agent that is ``learning''; when the agent's information influences the environment that is
``action'').

And similarly, now we have that reward is also, information. What is it information about?
Well, an RL agent produces probabilities $\pi(x)$ for its actions and samples from them (so
the policy also functions as the agent's beliefs about its own actions), and reward updates
these probabilities---so it is simply information about ``what I will do'', and is
definitionally rational, because ``what I will do'' is \emph{actually} updated and the
agent's beliefs about it are updated in-sync. You can specify \emph{any} objective
whatsoever, and optimizing it (at least through entropic mirror ascent) will still be
rational Bayesian inference. In other words, reward is \emph{hyperstitional} information.

(And if you multiply all the exponentiated rewards the agent gets through its training, you
get its final prior, which expresses its exponentiated utility function. Kinda, maybe,
oversimplified, whatever.)

In bandits, the objective is just the reward; in more general RL you must use Q-values, and
also do importance sampling (i.e.\ use $f(p)/p_i$ where $i$ is the taken action instead of
$f(p)$) because you only have access to the reward for one performed action. PPO methods
used in modern RL also are approximately entropic mirror ascent (TRPO is exactly it; other
stuff is kinda). You may also recall that evolution (replicator dynamics) is Bayesian
inference \citep[explained by][]{harperReplicatorEquationInference2009, baezInformationGeometryPart92012}. Indeed, it is exactly equivalent to this dynamic, with the fitness function as
an exponentiated reward.

\subsection{Continuous-time mirror ascent and information geometry}\label{agency:sec:infogeom}

It is interesting to take the continuous limit of this dynamics. We have:
$$\frac{d}{dt}\nabla G(x)= \nabla f(x)$$

(i.e.\ the velocity in the mirror space is the gradient of your objective function). What
does this look like in the original primal space? Well we can apply the chain rule to
rewrite the left side as $\nabla^2 G(x)\frac{dx}{dt}$ (where $\nabla^2$ is the Hessian)
thus:
\begin{equation}
\frac{dx}{dt}=[\nabla^2 G(x)]^{-1}\nabla f(x)\tag{$*\!*\!*$}
\end{equation}

This gives us a very natural interpretation of mirror ascent: if you let the Hessian
$\nabla^2 G(x)$ define a Riemannian metric on your space, \textbf{then mirror ascent is
simply gradient ascent of the objective $f(x)$ on that geometry.} Physicists are of course
quite used to transforming between primal and dual spaces by multiplying by the metric and
its inverse respectively.

Pitfall prevention: this is not gravity, nor is it any Newtonian force (as that would drive
the acceleration not the velocity).

When $G(x)$ is entropy, its Hessian is called the ``Fisher information matrix'' or
``Fisher--Rao metric'' $\diag(1/p)$ and the geometry defined by this Riemannian metric is
called information geometry.

\subsection{Advantage is rivalry}\label{agency:sec:advantage-rivalry}

Actually, the rule in the above section is not exactly right for the probabilistic case,
because it does not account for the probability normalization (i.e.\ the fact that the
original optimization itself has a constraint of probabilities adding up to 1).

You can just take the continuous limit of the multiplicative update rule directly.

In Bayesian inference, imagine the likelihood function to arrive over a period of time, so
that over time $dt$, you get infinitesimal evidence $\exp(v\ dt)$. Call $v$ the ``rate of
evidence'' (and it is a vector where each index is a hypothesis)---and $\bar{v}=p\cdot v$ is
your ``expected evidence rate''. Then (where $\odot$ is the pointwise product of vectors)
Bayes's theorem is:
$$p(t+dt)=\frac{p(t)\odot\exp(v\ dt)}{p(t)\cdot \exp(v\ dt)}$$

Take first-order approximations $\exp(x)\approx 1 + x$ and $\frac1{1+x}\approx 1-x$:
\begin{align*}
p(t+dt)&=\frac{p(t)\odot (1+ v\ dt)}{p(t)\cdot (1+ v\ dt)}\\
&=p(t)\odot \frac{1+ v\ dt}{1 + \bar{v}\ dt}\\
&= p(t)\odot (1+v\ dt) \odot (1-\bar{v}\ dt)\\
&= p(t)\odot (1+v\ dt-\bar{v}\ dt)\\
\implies dp/dt&= p(t)\odot(v-\bar{v})
\end{align*}
This is the differential form of Bayes's theorem, and is the replicator equation from
evolution.

In the RL/entropic mirror ascent case, we have:
\begin{align*}
p(t+dt) &=\frac{p(t)\odot \exp\left(\nabla f(p(t))\ dt\right)}{p(t)\cdot \exp\left(\nabla f(p(t))\ dt\right)}\\
&= p(t)\odot (1 +\nabla f(p(t))\ dt-p(t)\cdot \nabla f(p(t))\ dt)\\
dp/dt&= p(t)\odot (\nabla f(p)-p\cdot \nabla f(p))
\end{align*}
Which is the same replicator equation, with the gradient of the reward $\nabla f(p(t))$ as
the ``evidence rate'' and the expected reward $p(t)\cdot \nabla f(p(t))$ as the ``expected
evidence rate''.

But notice: this $\nabla f(p)-p\cdot \nabla f(p)$ is \textbf{Advantage}, the difference
between the $Q$ and $V$ values!

Without ``rivalry'' (the constraint that probabilities would add up to 1), you would just
have $dp/dt= p(t)\odot \nabla f(p)$ (equivalent to the equation from the previous
section---the inverse Fisher information is $\diag(p)$ so left-multiplying by it is
equivalent to doing $p(t)\odot$). \textbf{\emph{Advantage} arises from rivalry}, i.e.\ the
constraint that these probabilities must add up to 1 (indeed, using advantage instead of the
$Q$-value is only compulsory for tabular softmax parameterized policies).

\subsection{What I'm confused about: dynamics}\label{agency:sec:dynamics-confusion}

Finally, the thing I am actually confused about:

You know, a big problem in life is that people like expressing everything in terms of
equilibrium (in economics, game theory,
entropostatics aka thermodynamics \citep{liuEquilibriumThermoeconomics}) with little regard for the \emph{dynamics} that leads to the equilibrium.

(This is closely related to how we just model agents as rational (i.e.\ imagine that the
utility function can just be arg-maxed) instead of modeling bounded rationality which has
computational costs (and you can't just minus them from the objective either, because you
don't know what those costs are) and is fundamentally dynamic because you can't instantly
optimize something but that computational time is itself a cost\ldots)

Entropic mirror ascent seems like a very promising model of dynamics. Because evolution (the
most fundamental learning algorithm) actually follows it, and it is naturally a dynamical
Bayesian inference. And I have some work suggesting that economies can naturally be thought
of as doing it too.

But on the other hand \ldots\ I can't imagine how dynamical thermodynamics could have
anything to do with this? The simplest actually dynamical example in thermodynamics is
Newton's law of cooling, and that actually happens to be \textbf{gradient ascent on
entropy}, not continuous Bayesian inference/entropic mirror ascent.

Modeling thermodynamic dynamics fundamentally seems to have something to do with thinking
like ``Ok, I am not completely uncertain about this because I know where the system was 1
second ago; that is somehow a soft constraint'', i.e.\ a generalized notion of ``fraction''.
And I don't even know how gradient ascent on entropy comes out of this, maybe Newton's law
of cooling is just a coincidence?

% ============================================================
\section{Setup}\label{agency:sec:setup}
% Source: main article, Section 1.

Here is how we will define an ``economy''.

\begin{definition}[economy]\label{agency:def:economy}
There is a vector space $\Qspace$ in which quantity vectors live (so its dimensionality is
the number of unique goods). There is a directed graph, where every node $i$ is a
``factory'' labelled with production function $f_i:\Qspace\to\Qspace$ and every edge $ij$ is
labelled by a quantity vector $q_{ij}\in\Qspace$. The graph is subject to physical
constraints:
$$f_i(x_i)=y_i$$
where $x_i=\sum_k q_{ki}$ and $y_i=\sum_j q_{ij}$. Furthermore some nodes, called
``consumers'', have a utility function $U_i(x_i,y_i)$.
\end{definition}

We may imagine a global welfare objective $W = \sum_i U_i(x_i, y_i)$ (WLOG all the weights
are 1).

\textbf{U.S.O.W.} We will usually think of $(V,E)$ as a DAG and $f_i, U_i$ are continuous
and differentiable; utilities have no explicit dependence on flows or shares. Cyclic graphs
are certainly \emph{interesting}, but they are hard in general.

(\textbf{But what about factories choosing what to produce? Shouldn't $f_{i}(x_i)$ be a set
and $y_i$ be chosen from this set?} For now we will treat that as a set of every possible
$f_i$, and the choice of what to produce amounts to routing to the right $f_i$. Maybe this
is bad from a computational POV or whatever; alternatively we may consider the $f_i$'s
themselves as trainable sub-agents with some internal parameters.)

\subsection{General price recursion}\label{agency:sec:price-recursion}

Introduce Lagrange multipliers $\lambda_i$ for the production constraints:
$$\mathcal L = \sum_i U_i(x_i,y_i) -\sum_i \lambda_i\cdot (y_i-f_i(x_i))$$
Differentiate with respect to quantity vectors $q_{i\to j}$, we get the recursion:
\begin{equation}
\lambda_i=\nabla_{y_i}U_i+\nabla_{x_j}U_j+J_{f_j}(x_j)^\top\lambda_j\tag{1.1}\label{agency:eq:lambda-recursion}
\end{equation}
i.e.\ the price of output from $i$ equals the direct marginal utility/cost at $i$, plus the
direct marginal utility at $j$, plus the downstream production value obtained by sending the
good through $j$. The term $J_{f_j}(x_j)^\top\lambda_j$ is the backpropagation operation.

However, this does not directly let us perform backpropagation on the market graph, since
$q_{i\to j}$ are not free variables: they must satisfy the constraint
$\sum_j q_{i\to j}=y_i$. So instead parameterize quantity flows using routing shares:
$$q_{i\to j}=s_{ij}\odot y_i$$
$$s_{ij,g}\geq 0,\qquad \sum_j s_{ij,g}=1$$

The forward computation is (where $\odot$ indicates the elementwise product, and $r_i$ is an
exogenous resource endowment if any):
$$x_i = r_i+\sum_h s_{hi}\odot y_h$$
$$y_i=f_i(x_i)$$

Defining adjoints $a_i=\frac{\partial W}{\partial x_i}$ and
$b_i=\frac{\partial W}{\partial y_i}$, $c_{ij}=\frac{\partial W}{\partial s_{ij}}$ and
differentiating:
\begin{equation}
a_i=\nabla_{x_i}U_i+J_{f_i}(x_i)^\top b_i\tag{1.2}\label{agency:eq:adjoint-a}
\end{equation}
\begin{equation}
b_i=\nabla_{y_i}U_i+\sum_j s_{ij}\odot a_j\tag{1.3}\label{agency:eq:adjoint-b}
\end{equation}
\begin{equation}
c_{ij} = y_i\odot a_j\tag{1.4}\label{agency:eq:share-gradient}
\end{equation}
Now, these shadow price propagation equations themselves are just \emph{true}: not dependent
on any particular dynamics/optimization algorithm. They are highly \emph{suggestive} of a
backprop-like dynamics, and the fact that these prices actually are so important in
real-world economic decision-making suggests that this suggestiveness is right.

You can transform the shares into logit space to do normal backpropagation without simplex
constraints, or you can directly do optimization on a simplex, e.g.\ via entropic mirror
ascent which gives the update rule (with geometric learning rate $\eta$) as in
Section~\ref{agency:sec:cbi}. Regardless, the optimization rule is always: ``move outputs
towards buyers with higher marginal value $a_j$''. At equilibrium, all active destinations
of the same good from the same source have equal marginal value, i.e.\ for active edges
($q_{i\to j,g}>0$, $q_{i\to k,g}>0$):
$$a_{j,g}=a_{k,g}$$

% ============================================================
\section{Markets are MDPs; advantage = rivalrous goods}\label{agency:sec:mdp}
% Source: main article, Section 2.

Here is a nice interpretation of our economic model. Imagine there is only \emph{one} unit
good, i.e.\ a particle in the whole market. If the production functions are identity, then
its movement through the market is governed by policy $s_{ij,g}$, collecting rewards
$\nabla U$ along the way.

Production makes the particle ``branch'': at node $i$ a unit of good becomes a bundle
$J_{f_i} \cdot (\text{unit})$. We may regard the Jacobian $J_{f_i}$ as a ``matrix-valued
discount factor''. We also have a Bellman recursion:
\begin{equation}
\bar a_{i} := \sum_j s_{ij}\odot a_{j}\tag{2.1}\label{agency:eq:vbar}
\end{equation}
\begin{equation}
a_j = \underbrace{\nabla_{x_j} U_j+ J_{f_j}^\top\nabla_{y_j} U_j}_{\text{reward}} + J_{f_j}^\top \bar{a_i} \tag{2.2}\label{agency:eq:bellman}
\end{equation}

\begin{center}
\begin{tabular}{@{}ll@{}}
\toprule
RL & market \\
\midrule
state                         & ``one unit of good $g$ sits at node $i$''             \\
action                        & ``ship it to $j$''                                    \\
policy $\pi(j \mid i)$        & routing share $s_{ij,g}$                              \\
reward                        & marginal utility $\nabla U$ collected along the way   \\
action value $Q(i,j)$         & destination price $a_{j,g}$                           \\
state value $V(i)$            & average received price $\bar a_{i,g}$                 \\
advantage $A(i,j)$            & price premium $a_{j,g} - \bar a_{i,g}$                \\
visitation measure $d^\pi(i)$ & quantity present $y_{i,g}$                            \\
\bottomrule
\end{tabular}
\end{center}

Optimizing this market (computing the $s_{ij}$s) is then equivalent to solving this MDP. In
fact, we have a ``Policy Gradient Theorem'' (the $c_{ij}$
equation~\eqref{agency:eq:share-gradient} from earlier):
$$c_{ij} := \frac{\partial W}{\partial s_{ij}} = y_i\odot a_j$$

Or parameterizing with logits $s_{ij,g} = \exp(\theta_{ij,g}) / \sum_k \exp(\theta_{ik,g})$:
\begin{equation}
\frac{\partial W}{\partial \theta_{ij,g}}
  \;=\; s_{ij,g}\, y_{i,g} \left( a_{j,g} - \bar a_{i,g} \right)
  \;=\; q_{i\to j, g} \left( a_{j,g} - \bar a_{i,g} \right)\tag{2.3}\label{agency:eq:logit-gradient}
\end{equation}

\begin{proof}
Chain rule through the softmax Jacobian
$\partial s_{ik}/\partial \theta_{ij} = s_{ik}(\delta_{kj} - s_{ij})$ applied
to~\eqref{agency:eq:share-gradient}.
\end{proof}

In fact even the interpretation of $J_{f_i}^\top$ as a ``matrix-valued discount factor'' has
a natural interpretation. Generalizing $\gamma <1$, we have in a cyclic economy the
Hawkins--Simon condition $\rho(S^\top J_f) < 1$ where $S$ is the full matrix of $s_{ij}$'s
(for a particular good) and $\rho$ is the spectral radius (the largest absolute eigenvalue),
which tells us that a cyclic economy cannot be a perpetual motion machine.

(Well actually, it kinda can be a perpetual motion machine, if the dynamics of the cyclic
economy represent actual time rather than just convergence to an equilibrium. But tbh I
would rather represent that unrolled, as an infinite graph with nodes stretching out into
the future.)

(In the ``wide market RL'' I introduced in my paper
``Market-based architectures in RL and beyond'' \citep{pallavisudhirMarketbasedArchitectures2025}, I don't think of the MDP itself as
a market, but the transition probabilities are determined by an additional Garrabrant-like
program market of traders and the network is more ``fundamentally dynamic'' rather than
graph-like. Something in-between---dynamic like program markets but efficient like
graphs---like a graph that forms new connections or performs search of some sort when shadow
prices exceed some friction-induced background transaction cost, seems like a good model for
describing real markets, Knightian uncertainty etc.)

% ============================================================
\section{Chain markets are MLPs}\label{agency:sec:mlp}
% Source: main article, Section 3.

\draft{The source article invokes an ``alignment lemma'' from the research notes
accompanying the \texttt{marketgrad} library; we reproduce it here. Give each node $i$ the
objective (``utility-inclusive profit at posted prices $p$''):}
$$J_i \;=\; \underbrace{\sum_j p_j \cdot (s_{ij} \odot y_i)}_{\text{revenue}}
      \;-\; \underbrace{\sum_k \pi_i \cdot q_{k \to i}}_{\text{input bill}}
      \;+\; U_i(x_i, y_i),$$
\draft{where $p_j$ are prices posted by $i$'s buyers and $\pi_i$ is the price vector $i$
posts for its inputs; let node $i$'s own controls be its output shares $s_{i\cdot}$ and any
internal technology parameters $w_i$ of $f_i$.}

\begin{lemma}[alignment]\label{agency:lem:alignment}
If posted prices equal adjoints, $p_j = a_j$ for all buyers $j$ of $i$, then for every
control $u \in \{s_{i\cdot}, w_i\}$:
$$\nabla_u J_i \;=\; \nabla_u W .$$
\end{lemma}

\begin{proof}
For shares:
$\partial J_i / \partial s_{ij,g} = y_{i,g} p_{j,g} = y_{i,g} a_{j,g} = \partial W / \partial s_{ij,g}$
by~\eqref{agency:eq:share-gradient}; the input bill does not depend on $i$'s controls
(received quantities are set by upstream shares). For technology:
$\nabla_{w} J_i = (\partial y_i/\partial w)^\top (\nabla_{y_i} U_i + \sum_j s_{ij} \odot p_j) = (\partial y_i / \partial w)^\top b_i = \nabla_w W$.
\end{proof}

At shadow prices, \textbf{selfish first-order incentives coincide exactly with the welfare
gradient}---for routing decisions \emph{and} for technology investment. This is a
differential form of the first welfare theorem, and it is the economic content of backprop:
\emph{the price system is what aligns local greed with global credit assignment.}

\begin{theorem}[chain = MLP]\label{agency:thm:chainmlp}
Let the economy be a directed chain $0 \to 1 \to \cdots \to L$ with resource $r$ injected at
node $0$, production functions $f_\ell$, and a single consumer with utility $u$ at node $L$.
Then all shares are trivially equal to 1, and:
\begin{enumerate}
  \item The forward pass is the feedforward computation $h_{\ell+1} = f_\ell(h_\ell)$,
  $h_0 = r$, with welfare $W = u(h_L)$: \textbf{goods bundles are activation vectors} (one
  neuron per good; a node/factory is a whole layer).
  \item The market's shadow prices are the backprop gradients:
  $a_\ell = \partial W / \partial h_\ell = J_{f_\ell}^\top \cdots J_{f_{L-1}}^\top \nabla u$.
  \item \textbf{Internal parameterization.} If $f_\ell$ has parameters $w_\ell$, then by
  Lemma~\ref{agency:lem:alignment}, gradient ascent by each firm on its own profit at posted
  prices is identical to backprop/SGD on the network's weights:
  $\nabla_{w_\ell}(\text{profit}_\ell) = \nabla_{w_\ell} W$.
  \item Conversely, every MLP whose activations are nonnegative (i.e.\ with ReLU or similar
  activation functions) arises this way, with \textbf{bias vectors = exogenous endowments}
  $r_\ell$ injected at layer $\ell$ and the objective function as the terminal consumer's
  utility.
\end{enumerate}
\end{theorem}

\begin{proof}
1--2 are the forward/backward recursions of Section~\ref{agency:sec:setup} with a single
out-edge; 3 is Lemma~\ref{agency:lem:alignment}; 4 is construction.
\end{proof}

Theorem~\ref{agency:thm:chainmlp} completely formalizes the ``economies = backprop/neural
networks'' intuition that people have been groping at. Note that economies are more
\emph{general} than classical MLPs: they are neural networks with \emph{routing} (which is
kind of ``transposed static attention'': in markets, the goods are scarce while in
transformers, the receiver-side attention budget is scarce).


% ============================================================
\section{Continuous Bayesian inference as a dynamics for markets}\label{agency:sec:cbi}
% Source: main article, Section 4.

In Section~\ref{agency:sec:hyperstition}
(``Reward as hyperstitional information'', \citealt{pallavisudhirRewardHyperstitional2026}), I conveyed that entropic mirror ascent (aka continuous-time
Bayesian inference) is the natural dynamics for probabilities. Rewards and such are just an
``evidence rate''.

It seems very suggestive that our routing shares are literally just probabilities (of ``a
particle of good being sent to some factory''). Indeed, we can let the dynamics of our
market be ``entropic mirror ascent on total welfare'' i.e.\ update the shares each time step
as (indices $i,j,g$ are implicit):
$$s'=\argmax_{s'} \{\eta\langle \partial W/\partial s, s'\rangle - \KL(s'\|s)\}$$
From~\eqref{agency:eq:share-gradient} we know that the share gradient equals the revenue
rate $\partial W/\partial s = ya =: c$ thus the solution to this is (read
Section~\ref{agency:sec:hyperstition} if this isn't immediately clear):
\begin{equation}
s' \propto s \exp{(\eta c)}\tag{4.1}\label{agency:eq:mwu}
\end{equation}

Differentiate with respect to time to get the usual replicator flow/continuous-time Bayesian
inference:
\begin{equation}
\dot s = s\left( c - \bar c \right) \tag{4.2}\label{agency:eq:replicator}
\end{equation}

So: \textbf{the market (under entropic mirror ascent) is a Bayesian learner}, with each
agent doing ``hyperstitious'' learning. I.e.\ each agent $i$ has a belief distribution
$s_{i,-}$ on the question ``where will I send an infinitesimal quantity to?'' and the
infinitesimal evidence it receives for each hypothesis $j$ is exponentiated revenue ratio
$\exp{(\eta c_{ij})}/Z$.

\begin{thesisclaim}
Revenues are rewards are information (log-likelihoods).
\end{thesisclaim}

You can use this equivalence to transfer any result.

\begin{center}
\small
\begin{tabular}{@{}llp{4.2cm}p{5.2cm}@{}}
\toprule
domain & mass & fitness / evidence & equilibrium \\
\midrule
evolution & population shares      & reproductive fitness & fitness equalized (Fisher) \\
markets   & supply shares / wealth & prices, revenue      & law of one price, no arbitrage \\
inference & posterior probability  & log-likelihood       & likelihoods equalized/posterior support \\
RL        & policy $\pi(j \mid i)$ & $Q$-values ($d^\pi(i)\,Q(i,j)$ visitation-weighted; advantage $A = Q - V$ as the relative form) & advantage vanishes on support, $\le 0$ off it---Bellman optimality/policy improvement stops; generally a greedy deterministic policy, or Boltzmann $\pi \propto e^{Q/\tau}$ under entropy regularization \\
\bottomrule
\end{tabular}
\end{center}

Much like Fisher's fundamental theorem of natural selection we have the ``fundamental
theorem of welfare dynamics'': fix the technologies $f_i$ and utilities $U_i$ and let all
shares evolve simultaneously by~\eqref{agency:eq:replicator}. Then
\begin{equation}
\dot{W}
  = \sum_{i \in V} \sum_{g \in G} y_{i,g}^2\;
        \Var_{j \sim s_{i\cdot,g}}\!\left( a_{j,g} \right) \tag{4.3}\label{agency:eq:ftwd}
\end{equation}

\begin{proof}
$\dot W = \sum_{i,j,g} \frac{\partial W}{\partial s_{ij,g}} \dot s_{ij,g}$ (since the
partials are total derivatives). Substituting~\eqref{agency:eq:share-gradient}
and~\eqref{agency:eq:replicator}, each $(i,g)$ block contributes
$y_{i,g}^2 \sum_j s_{ij,g} a_{j,g}(a_{j,g} - \bar a_{i,g}) = y_{i,g}^2 (\mathbb{E}[a^2] - \bar a^2)$.
\end{proof}

% ============================================================
\section{Markets rationalize reward and identity}\label{agency:sec:rationalize}
% Source: main article, Section 5.

I started my MATS work with the following observation:

\begin{quote}
Utility functions and hyperstitional ``identities'' (in the sense of predictive
processing---i.e.\ you do what you predict you'll do) are the same thing.
\begin{enumerate}
  \item Both are ``arbitrary'': the choice of utility function, nor the co-ordinated choice
  of equilibrium (or whatever) are determined by existing rules of rationality.
  \item Purely predictive agents (i.e.\ \textbf{non-embedded} agents, \emph{not}
  Predict-o-matic) do not have identities; their beliefs are completely specified by Bayes
  (in the limit of all information being supplied). [I mean sure you could say that some
  priors may not allow convergence to the same thing, or debate what it means to take the
  limit of ``all information being supplied'', or talk about boundedly rational agents
  having ``globally inconsistent'' beliefs, but this seems like a qualitatively different
  thing.]
\end{enumerate}
So ideally, we would like to have a theory where internal sub-agent co-ordination is
equivalent to a choice of utility function. More precisely: the theory should connect this
to \emph{reward} (which ultimately ``selects'' or trains the identity), and build a
satisfying picture that explains how that exogenous reward itself fits into a broader
predictive-processing-only picture.
\end{quote}

Or more simply: rational behaviour for a \emph{purely predictive}/non-embedded agent is
completely specified by Bayes's theorem, but for any agent that performs \emph{actions} in
the world/whose beliefs affect the world, there is an arbitrary thing called the ``utility
function''. The utility function is shaped by reward, but then reward is arbitrary.

I continued:

\begin{quote}
The picture I expect right now is vaguely something like:
\begin{itemize}
  \item ``the universe as a whole'' (or closed systems more generally) do not have
  identities/utility functions, but individual sub-agents do
  \item Identities are the result of the universe hierarchically enforcing its will on its
  sub-agents via reward/selection. [But the question of exactly how actions i.e.\ things
  causally determined by agent beliefs, emerge, should be clarified].
  \item Reward has to do with constraints in statistical mechanics
\end{itemize}
This intuition is largely based on Point 2, and also the suggestive connection between
evolutionary fitness functions and Bayesian evidence
\citep[``The Replicator Equation as an Inference Dynamic'',][]{harperReplicatorEquationInference2009}, which supports the idea that evolutionary selection is somehow the universe
``inserting its information'' into an agent.
\end{quote}

Now observe how markets have two processes going on simultaneously:
\begin{enumerate}
  \item \textbf{Agents learn the shares $s_{ij}$}, e.g.\ through continuous-time Bayesian
  inference, with prices (or rather revenue) as the ``information'' or reward supplied.
  \item \textbf{Price recursion} i.e.\
  equations~\eqref{agency:eq:lambda-recursion}--\eqref{agency:eq:share-gradient} or
  equivalently \eqref{agency:eq:vbar}--\eqref{agency:eq:bellman} which follow a
  Bellman/backprop-like rule.
\end{enumerate}

\begin{quote}
\emph{FAQ: Wait, how can price recursion be equivalent to both Bellman and backprop?}

Backprop/reverse-mode autodiff on \emph{any} computation graph produces a ``Bellman-shaped''
recursion, i.e.\ value of a node = sum of its consumer nodes' values transformed by
$J_{f_i}^\top$.

What makes the market price recursion actual Bellman (with matrix-valued discount factors
$J_{f_i}^\top$) rather than merely ``Bellman-shaped'' is the simplex constraint which lets
the routing shares be interpreted as probabilities.
\end{quote}

In other words:

\begin{thesisclaim}
The economy is exactly the ``super-agent'' that generates reward to its sub-agents, and
price recursion is the ``rational theory of reward''.
\end{thesisclaim}

In Section~\ref{agency:sec:hyperstition}
(``Reward as hyperstitional information'', \citealt{pallavisudhirRewardHyperstitional2026}), reward was entirely an underdetermined hyperstition.
Markets (equations \eqref{agency:eq:vbar}--\eqref{agency:eq:bellman}) tell you how to
determine it.

\draft{The main article's appendix of caveats and open questions is reproduced in
Appendix~\ref{agency:sec:caveats}.}

% ============================================================
\section{The \texttt{marketgrad} library}\label{agency:sec:library}
% Source: marketgrad/README.md and marketgrad/SPEC.md (Sections 1 and 21).

\code{marketgrad} is a NumPy-first library for market-like optimization on production
graphs. It provides graph primitives, differentiable production and utility functions,
shares-based welfare optimization, cyclic fixed-point and adjoint solvers, simple
decentralized market dynamics, primal-dual optimization scaffolding, diagnostics,
visualizations, examples, and tests. \draft{The theory of this chapter is maintained
alongside the library: its identities are checked numerically against the implementation
in the library's test suite.}

The core abstraction is a directed graph of productive, consumptive, and resource nodes.
Goods flow along edges. Nodes transform input bundles into output bundles via production
functions. Some nodes receive utility from input and/or output bundles. The library supports
several algorithmic interpretations of such systems:
\begin{enumerate}
  \item \textbf{Shares-based welfare optimization.} A centralized, physically feasible,
  backpropagation-like optimizer where edge quantities are parameterized by routing shares.
  \item \textbf{Decentralized market-like dynamics.} A more realistic simulation of agents,
  bids, prices, wealth, budgets, ownership, sales, and production over time.
  \item \textbf{Primal-dual constrained optimization.} A general
  Lagrangian/augmented-Lagrangian/ADMM-style solver that can handle constraints not easily
  captured by shares, possibly allowing temporary physical infeasibility.
\end{enumerate}

The central conceptual bridge is:
\begin{center}
\code{backprop adjoints $\approx$ shadow prices $\approx$ Lagrange multipliers}
\end{center}
and:
\begin{center}
\code{gradient ascent over feasible routing decisions $\approx$ reallocating goods toward
higher marginal value}
\end{center}

The library supports both acyclic and cyclic production graphs. In acyclic graphs,
quantities can be computed by a topological forward pass and prices by a reverse backward
pass. In cyclic graphs, quantities are defined by fixed-point equations and prices by
implicit adjoint equations, as in deep equilibrium models.

The primary clean algorithm is shares-based optimization:
\begin{quote}
forward pass: compute feasible quantities \\
backward pass: compute adjoint prices \\
update shares: route goods toward higher marginal value
\end{quote}
For acyclic graphs, this is ordinary forward/backward propagation. For cyclic graphs, this
becomes a fixed-point solve for quantities and an implicit adjoint solve for prices. The
default share update is entropic mirror ascent,
\code{s\_new $\propto$ s\_old * exp(learning\_rate * marginal\_value)}; projected gradient
and softmax-logit optimization are also supported, but as alternative update geometries. The
decentralized market simulator is the realistic economic dynamics layer (agents bid, trade,
produce, earn, spend, and update prices locally), and the primal-dual solver is the advanced
constrained-optimization backend (quantities and multipliers co-evolve, possibly with
temporary constraint violations).

The library's purpose is to make the analogy between markets and backpropagation precise,
programmable, visualizable, and experimentally testable. JAX and PyTorch backends are
optional (\code{marketgrad[jax]}, \code{marketgrad[torch]}); without either extra, the
default NumPy backend and finite-difference autodiff remain unchanged.

%% TODO(agency): SPEC.md is a full library specification (~2400 lines: data structures,
%% solver families, market-tick protocol, primal-dual/ADMM backends, visualization and
%% testing requirements); only its purpose section and summary are lifted here.

% ============================================================
% Subappendices
% ============================================================
\begin{subappendices}

\section{Caveats and questions}\label{agency:sec:caveats}
% Source: main article, appendix.

\ldots well, with two caveats:

\textbf{I. Our market model still has utility functions at the end of the tunnel.} Sure you
can probably replace these with probabilities active inference style, but the point is they
are still under-determined. So markets tell you how a super-agent gives rewards to its
sub-agents, and maybe that super-agent gets rewards from its super-agent etc.\ but what is
there ``all the way up''?

E.g.\ factories' rewards are derived from human consumers' utility functions, humans'
rewards are determined by evolution's fitness functions, evolution's fitness functions are
themselves hyperstitional but maybe there is a grand cosmic-scale selection process that
selects good evolutionary priors over bad evolutionary priors, and maybe everything somehow
boils down to ``the universe wants to maximize its entropy'' or something like that.

Thermodynamics (free entropies) seems relevant to this.

\textbf{II. Markets are just a model, right?} It's not the case that \emph{all} reward is
computed based on price recursion (in fact nothing does perfectly so); rather we generally
use heuristics. But maybe this is like how nobody actually follows Bayesian reasoning, but
Bayes is still the ideal or optimal mechanism in some sense.

This would need a Dutch book type result.

And a third question, not crucial to the ``price recursion rationalizes reward'' philosophy
but the continuous Bayesian inference dynamic which I still find important.
%% TODO(agency): the source sentence trails off here ("...which I still find important
%% because .").

\textbf{III. Wealths and perfect competition.} We kinda just postulated continuous Bayesian
inference as the dynamics of markets because it feels right. A proper justification of this
would require a game-theoretic consideration of agents maximizing their local objectives
(presumably with an Aumann-style agent continuum) and having wealths/budgets.

AI tells me continuous Bayesian inference has a natural interpretation in terms of wealths
(which kinda makes sense; accumulated evidence is wealth after all right).

Another question I am interested in, that I feel is related but am not sure:

\textbf{IV. ``Dynamic rationality''.}

It seems to me that dynamics is fundamental to understanding ``bounded rationality'': the
fact that you cannot instantly argmax a utility function or compute a market equilibrium.
The fact that you can have systems that are ``dynamically complete'' (there is eventually an
Arrow--Debreu security for every event in the world) and ``dynamically consistent'' (at most
a finite amount of arbitrage can be extracted from the market in infinite time) but not
complete and consistent (G\"odel's theorem) also seems suggestive.

Questions:
\begin{enumerate}
  \item Continuous Bayes feels like a natural model of dynamics: only an infinitesimal
  amount of information gets processed in infinitesimal time. How does this relate to all
  the other things about bounded rationality, like logical non-omniscience and Knightian
  uncertainty?
  \item But Continuous Bayes only models \emph{inserting} information, what about
  \emph{relaxing} constraints/information (e.g.\ Newton's law of cooling)? Is there a model
  of the latter, and is it important? Feels like it might be relevant to Question I, where I
  spoke of thermodynamics.
  \item It seems as though dynamics too may be ``arbitrary'' like reward, and like reward,
  ``selected by the super-agent''---i.e.\ markets prefer agents who learn, they prefer
  agents who learn faster, respond faster to shocks, etc.
\end{enumerate}

\end{subappendices}
